\section{Introduction}
There are two distinct realization questions. A terminal numerical
interface asks for a correct mean decoder after a word. A causal
observation experiment asks for one consistent law of all outputs,
including controls selected from earlier outputs. With repeatable
probes, Section~\ref{sec:causal59} proves that bounded width at any
fixed total transcript error below one requires a finite future-response
quotient, and that this quotient itself gives an exact realization.
This conclusion is stronger than a radius threshold for numerical
marginals and has a different proof. The terminal theory is developed
in full thereafter.

We ask whether a free clock and a separate stochastic realization at
every horizon can reduce the limiting number of hidden labels needed to
preserve a continuous numerical interface. Under eventual approximate
returns, they cannot: the least clocked width converges to the stationary
minimum at the same closed error, and finite limits are eventually
attained. More strongly, arbitrarily many cuts of a fixed width already
force a stationary realization of that width, even when all other
registers are unrestricted.

For compact physical groups, stationary transition randomness and hidden
permutations above the physical identity can be removed without losing
a label or increasing error. Randomized initialization remains. This
identifies the stationary minimum with a finite continuous physical
action and yields an equality-inclusive finite-quotient criterion,
including groups with infinitely many components. A change of frame
identifies each phase minimum without a phase-register overhead.

The command-width resource immediately below excludes any sampled answer
register. The full proofs of the structural reductions, component-radius
localization, phase and profinite refinements, algebraic effectivity, and
explicit finite-group complexity are contained in this article. The
separate quantitative paper concerns orbit geometry and an explicit
rational exact--noisy separation; none of its theorems is a premise here.
\subsection{The interface and its two resources}
Unless a section explicitly specifies a different physical space, $H$ is a compact Hausdorff group and $\mathcal A$ is a finite nonempty alphabet whose letters $u_a\in H$ generate a dense subgroup. The execution convention is
\[
                  u_{a_1\cdots a_n}=u_{a_n}\cdots u_{a_1}.
\]
There are finite nonempty seed and query sets $\mathcal S,\mathcal J$ and continuous maps $F_{s,j}:H\to Y_j$. In the structural results $Y_j$ is a nonempty norm-compact convex subset of a real Banach space. The localization and algebraic sections specify finite-dimensional output spaces. For a categorical law, $Y_j=\Delta_{M_j}$ and the norm on differences is $\frac12\|\cdot\|_1$, namely total variation. A prescribed joint law is one query; incompatible measurements are not assigned an artificial joint law.

\begin{definition}[Stationary realization]\label{def:stationary55}
A realization on $k$ labels consists of initialization rows $\alpha_s\in\Delta_k$, row-stochastic matrices $T_a$ of order $k$, and decoder values $D_j(i)\in Y_j$, all independent of time and word length. Its error is
\[
       \sup_{s,j,w\in\mathcal A^*}
        \|\alpha_sT_{a_1}\cdots T_{a_n}D_j-F_{s,j}(u_w)\|_j.
\]
The empty word is included. Let $M_\varepsilon$ be the least $k$ at error at most $\varepsilon$, or $\infty$ when none exists. A permutation realization has permutation command matrices; its initialization need not be deterministic.
\end{definition}
The clocked resource $W_{N,\varepsilon}$ instead minimizes the maximum size of the registers $S_0,\ldots,S_N$ at a fixed horizon, with arbitrary row-stochastic transitions depending on time and horizon, and a terminal query decoder. Definition~\ref{def:machine54} states the model formally. Every advertised persistent label counts, including unused padding. Private information retained for future commands is part of the current label. The original width model does not charge the external clock, real table descriptions, their construction, exact arithmetic, or exact atomic sampling. Section~\ref{sec:complexity56} gives separate finite-description and fair-random-bit bounds; it does not reinterpret those resources as free computational memory. A numerical decoder vector is not an additional sampled-answer register.

For an open normal subgroup $U\le H$ put
\begin{equation}\label{eq:quotientradius55}
 r(U)=\max_{s,j,h\in H}\rad_{Y_j}(F_{s,j}(hU)),
 \qquad\rad_Y(A)=\min_{y\in Y}\sup_{x\in A}\|x-y\|.
\end{equation}
The coset images are compact by continuity. The radius depends continuously on the translate, and compactness gives the displayed extrema. The restriction of the center to the legal output set is essential. The component radius $R_0$ uses $U=H^\circ$, whether or not this subgroup is open; it is the infimum of the finite-quotient radii. At the critical error, the relevant issue is attainment of that infimum, not a limit from larger errors.


The principal dependencies are finite stationary witnesses and Ramsey
extraction for clock removal, classical recurrent-simplex structure for
purification, and physical-kernel averaging for descent. The associated
new quantifiers are the same width, the same closed numerical error,
randomized initialization and arbitrary intervening register widths.
The classical semigroup statements themselves are not claimed as new.
